\documentclass[a4paper,11pt]{article}

\usepackage[ngerman]{babel}
\usepackage[top=3cm,bottom=3cm,left=3cm,right=3cm]{geometry} %NICHT IN KEMIE2


%\usepackage{microtype} %schöneres typesetting  mit [stretch=10] für nur 1% anstelle von 2%
%\usepackage{lmodern} %Latin Modern FONT
%\usepackage{pifont} %schönere Font



\usepackage{amsmath}
\usepackage{nicefrac}
\usepackage{amssymb}
%\usepackage{mathtools}

\usepackage{titlesec} %Um Section, etc bearbeiten zu können
\usepackage{setspace}
\usepackage{enumitem} %enumerate
\usepackage{multicol}
\usepackage[many]{tcolorbox}
\usepackage{hyperref}
\usepackage{arydshln}

%\renewcommand{\ttdefault}{lmtt} %Schriftart wird geändert zu Latin Modern Typewriter


\hypersetup{hidelinks,
			colorlinks=true,
			linkcolor=black,
			citecolor=miudiutex,
			urlcolor=miugray2
			}

\setlength{\parindent}{0pt}
\setcounter{section}{0}


\titlespacing*{\section}{0pt}{20mm}{6mm}
\titlespacing*{\subsection}{0pt}{6mm}{3mm}
\titlespacing*{\subsubsection}{0pt}{3mm}{0mm}



%\definecolor{miudiutex}{HTML}{2f3d39}
\definecolor{miudiutex}{HTML}{1d2926}
\definecolor{miugray}{HTML}{b4cccf}
\definecolor{miugray2}{HTML}{4d7365}
\definecolor{red}{HTML}{cf1f5c}
\definecolor{green}{HTML}{00A36C}
\definecolor{delphinidin}{HTML}{315794}
\definecolor{pelargonidin}{HTML}{b33807}
\definecolor{cyanidin}{HTML}{ba0746}
\definecolor{petunidin}{HTML}{b56fed}



\raggedcolumns %wichtig für Multicolumns


%================================
% SYMBOLS
%================================

\newcommand{\RA}{$\rightarrow$~}
\newcommand{\RL}{$\rightleftharpoons$~}
\newcommand{\HARP}{\ding{226}~}
%$\xrightarrow{2}$ %Pfeil mit 2 drüber


%================================
% SHORT COMMANDS (BOXES ; ETC)
%================================


\newcommand{\notiz}[1]{\begin{notizz}{}{}\begin{center}
			{\color{miugray2!50!miudiutex}#1} \end{center}\end{notizz}}
\newcommand{\zit}[2]{\begin{zitat*}{#1}{}\small #2 \end{zitat*}}

\newcommand{\frage}[2]{\begin{qst*}{#1}{}\small #2 \end{qst*}}

\newcommand{\unten}{\end{center} \tcblower \begin{center}}


\newcommand*\circled[1]{\tikz[baseline=(char.base)]{
            \node[shape=circle,draw,inner sep=0.5pt,miudiutex] (char) {\tiny #1};}}
            
\newcommand{\miusquare}{{\color{miugray2!50}\scriptsize $\blacksquare$}~}
\newcommand{\miusquarp}{{\color{petunidin!50}\scriptsize $\blacktriangleright$}~}
\newcommand{\niu}{\item[\miusquare]}
\newcommand{\nip}{\item[\miusquarp]}

%%%%% ABSAETZE

\newcommand{\pps}{\par \vspace*{1mm}}
\newcommand{\pp}{\par \vspace*{3mm}}
\newcommand{\ppbig}{\par \vspace*{8mm}}
\newcommand{\pphuge}{\par \vspace*{10mm}}

%================================
% FRAGE BOX
%================================

\tcbuselibrary{theorems,skins,hooks}
\newtcbtheorem{qst}{Frage:}
{%
	enhanced,
	breakable,
	colback = gray!15!white,
	frame hidden,
	boxrule = 0sp,
	borderline west = {2.5pt}{0pt}{red!70},
	sharp corners,
	detach title,
	before upper = \tcbtitle\par\smallskip,
	coltitle = black,
	fonttitle = \bfseries\sffamily,
	description font = \mdseries,
	separator sign none,
	segmentation style={solid, gray!50!black},
}
{th}




%================================
% ZITAT BOX
%================================

\tcbuselibrary{theorems,skins,hooks}
\newtcbtheorem{zitat}{Zitat}
{%
	enhanced,
	breakable,
	colback = gray!15!white,
	frame hidden,
	boxrule = 0sp,
	borderline west = {2.5pt}{0pt}{black!70},
	sharp corners,
	detach title,
%	before upper = \tcbtitle\par\smallskip,
	coltitle = gray!50!black,
	fonttitle = \bfseries\sffamily,
	description font = \mdseries,
%	separator sign none,
	segmentation style={solid, gray!50!black},
}
{th}



%================================
% SIMPLE SCHWARZER RAHMEN BOX
%================================


\newcommand{\boxfr}[1]{
	\begin{tcolorbox}[
	drop shadow southeast,
	enhanced,
	colframe=black!50,
	colback=white,
	boxrule = 1pt, 
%	toprule=0pt, bottomrule=0pt,rightrule=0pt,leftrule=0pt,
	boxsep=5pt,
	arc=1pt,
	]
	\begin{center}
	#1
	\end{center}
	\end{tcolorbox}
}

\definecolor{miudiutex}{HTML}{1d2926}
\definecolor{miugray}{HTML}{b4cccf}
\definecolor{miugray2}{HTML}{4d7365}
\definecolor{white2}{HTML}{f5f7f8}
\definecolor{red}{HTML}{cf1f5c}
\definecolor{green}{HTML}{00A36C}
\definecolor{delphinidin}{HTML}{315794}
\definecolor{uniblau}{HTML}{004e9f}
\definecolor{unigelb}{HTML}{fcba00}
\definecolor{unigrau}{HTML}{909085}
\definecolor{pelargonidin}{HTML}{b33807}
\definecolor{cyanidin}{HTML}{ba0746}
\definecolor{paeonidin}{HTML}{d15ca6}
\definecolor{petunidin}{HTML}{b56fed}
\definecolor{malvidin}{HTML}{8a2d8c}

\newcommand{\metermilli}{\mskip3mu \text{mm}}
\newcommand{\cm}{\mskip3mu \text{cm}}
\newcommand{\m}{ \text{m}}
\newcommand{\lite}{ \text{l}}
\newcommand{\kg}{ \text{kg}}
\newcommand{\g}{ \text{g}}
\newcommand{\km}{ \text{km}}
\newcommand{\h}{ \text{h}}
\newcommand{\s}{ \text{s}}
\newcommand{\tonne}{ \text{t}}
\usepackage[utf8]{inputenc}
\usepackage[T1]{fontenc}
\usepackage{lmodern}

\usepackage[
    activate={true,nocompatibility},
    final,
    tracking=true,
    kerning=true,
    spacing=true,
    factor=1100,
    stretch=30,
    shrink=20
]{microtype}
\usepackage{pdfpages}
\usepackage{awesomebox}
\usepackage{marginnote}
\tcbuselibrary{documentation}
\usepackage{sidenotes}
\usepackage{import}
\usepackage{xifthen}
\usepackage{transparent}
\usepackage[table]{xcolor}


\newcommand{\abbicon}{%
  \protect\tikz[baseline=0.2mm,scale=0.8]{%
    % Das blaue Quadrat mit abgerundeten Ecken
    \fill[uniblau, rounded corners=1.2pt] (0,0) rectangle (0.8em, 0.8em);
    % Der weiße Kreis in der Mitte
    \fill[white] (0.4em, 0.4em) circle (0.12em);
  }%
} 
\usepackage[labelfont={bf},font={color=miudiutex,small},figurename={\abbicon $~$Abbildung}]{caption}


\newcommand{\bckgrnd}{
\AddToHookNext{shipout/background}{
\begin{tikzpicture}[remember picture, overlay]
 \draw[draw=none,fill=uniblau!70, shift={(current page.north west)}] (-3,0) rectangle (23,-3.75);
 \draw[draw=none,fill=miugray, shift={(current page.north east)}] (-7,0) rectangle (3,-3.75);
\end{tikzpicture}
}
}



\newcommand{\incfig}[2]{%
\def\svgwidth{#2\columnwidth}
\import{C://LaTeX-Files/pngs/}{#1.pdf_tex}
}

\renewcommand{\textbf}[1]{\textsf{{\bfseries {#1}}}}



\titleformat{\section}[hang]
		{\LARGE \color{uniblau!75!black} \sffamily \bfseries}
		{\thesection}
		{6mm}
		{}
		[]
\titleformat{\subsection}[hang]
		{\large \sffamily \bfseries \color{black}}
		{\thesubsection}
		{4mm}
		{{\color{miugray}$\blacksquare ~$}}

\titleformat{\subsubsection}[block]
		{\normalsize \bfseries \sffamily \color{miudiutex}}%
		{\normalsize \color{black} \thesubsubsection}
		{2mm}
		{{\color{miugray}$\blacksquare ~$}}
		[ \color{black}]
		
\titleformat{\paragraph}
		{\normalsize \bfseries \sffamily \color{black}}
		{\footnotesize \theparagraph}
		{1mm}
		{}

\pagestyle{empty}
\titlespacing*{\section}{0pt}{20mm}{12mm}
\titlespacing{\section}{0pt}{20mm}{12mm}

\titlespacing*{\subsection}{0pt}{14mm}{4mm}
\titlespacing{\subsection}{0pt}{14mm}{4mm}

\titlespacing*{\subsubsection}{0pt}{6mm}{3mm}
\titlespacing{\subsubsection}{0pt}{6mm}{3mm}

\titlespacing*{\paragraph}{0pt}{6mm}{2mm}
\titlespacing{\paragraph}{0pt}{6mm}{2mm}


\let\oldmarginfigure\marginfigure
\let\endoldmarginfigure\endmarginfigure

\let\oldmarginfigure\marginfigure
\let\endoldmarginfigure\endmarginfigure

\renewenvironment{marginfigure}[1][0pt] % [1] steht für 1 Argument, [0pt] ist der Standardwert
  {\oldmarginfigure[#1]\captionsetup{labelfont={sf,bf,color=uniblau!75!black},font={color=miudiutex,small}}}
  {\endoldmarginfigure}
  
  
  
  
\renewcommand{\labelitemi}{\color{uniblau!70}$\bullet$}
\renewcommand{\labelitemii}{\color{uniblau!70}$\cdot$}
\newcommand{\plmtsquare}{{\color{uniblau!70}$\blacksquare~$}}
\newcommand{\splmt}[1]{{\color{uniblau!75!black} \sffamily #1}}
\newcommand{\fplmt}[1]{{\color{uniblau!75!black} \sffamily \bfseries #1}}

\newcommand{\antwort}[1]{
\begin{tcolorbox}[
	enhanced,
	enforce breakable,
	standard jigsaw,
	boxrule=0.7pt,
	colframe=black,
	sharp corners,
	boxsep=5pt,
	title=\bfseries \sffamily Antwort,
	opacityback=0,
	coltitle=miudiutex,
	attach boxed title to top text left={yshift=-\tcboxedtitleheight/2},
	boxed title style={colback=white,colframe=white},
	bottomrule at break=0pt,
	toprule at break=0pt,
	pad at break=16mm,
]
	\begin{center}
	#1
	\end{center}
	\end{tcolorbox}
}




\rowcolors{2}{miugray!50}{miugray!50}
\arrayrulecolor{white} % Weiße Gitterlinien
\setlength{\arrayrulewidth}{1.5pt} % Etwas dickere Linien
\renewcommand{\arraystretch}{1.4} % Mehr Zeilenhöhe für bessere Lesbarkeit}


%\setlist[itemize]{itemsep=0mm}
\setlist[enumerate]{itemsep=0mm}

\usetikzlibrary{arrows.meta, decorations.pathmorphing}


\newif\ifshowme
\showmetrue




\begin{document}
\ifshowme
\newgeometry{top=1.8cm,bottom=3cm,left=2cm,right=4cm}
\else
\newgeometry{top=1.8cm,bottom=4cm,left=1.5cm,right=7.5cm,marginparsep=-2.00cm,marginparwidth=4.75cm}
\fi

\bckgrnd


\begin{center}
\LARGE
\color{white}
\textbf{Tafelanschrift} \par \vspace*{2mm}
\large \color{miugray} \textbf{Prozessbezogene Lebensmitteltechnologie}
\end{center}
\par 
\vspace*{6mm}


\color{miudiutex}
\section*{Trocknung}
Gegenstromprinzip bei hygroskopischen Gütern oder temperaturunempfindlichen Gütern, die eine schnelle Trocknung erfordern.
\pp
Agglomerate sind besser löslich in Heißgetränken oder ähnlichem 
\subsection*{Gefriertrockner}
Sublimatiton 1kg Wasser entspricht 98.000 m$^3$ Dampf

\subsection*{Dipolrotation}


\subsection*{Konventionelle Erhitzung}
\begin{figure}[h]
\centering
\begin{tikzpicture}[
    >=stealth,
    very thick,
    color=blue!40!black, % Angelehnt an die dunkelblaue Tinte
    font=\sffamily\Large
]

% ==============================
% 1. Ofen
% ==============================
\node at (1.65, 3.5) [anchor=south west, inner sep=0] {Ofen};
\draw (0,0) rectangle (4.5, 3);

% Das Erhitzungsgut (Oval)
\draw (2.2, 1.5) ellipse (0.8cm and 0.6cm);

% Wärmestrahlung von links (geschwungene Pfeile)
\draw[->, decorate, decoration={snake, amplitude=1.5mm, segment length=5mm, post length=2mm}] (0.2, 1.2) -- (1.3, 1.2);
\draw[->, decorate, decoration={snake, amplitude=1.5mm, segment length=5mm, post length=2mm}] (0.2, 0.5) -- (1.5, 0.8);

\draw[->, decorate, decoration={snake, amplitude=1.5mm, segment length=5mm, post length=2mm}] (4.3, 2) -- (3.1, 1.7);
\draw[->, decorate, decoration={snake, amplitude=1.5mm, segment length=5mm, post length=2mm}] (4.3, 1.3) -- (3.1, 1);
%
%% Luftströmung/Strahlung von rechts (gebogene Pfeile)
%\draw[->] (4.3, 1.3) to[out=180,in=270] (3.5, 1.5);
%\draw[->] (4.3, 2.0) to[out=180,in=90] (3.5, 1.8);

% Pfeile von oben und unten
\draw[->] (1.8, 3.0) -- (1.8, 2.2);
\draw[->] (2.2, 3.0) -- (2.2, 2.2);
\draw[->] (2.6, 3.0) -- (2.6, 2.2);

\draw[->] (1.8, 0.0) -- (1.8, 0.8);
\draw[->] (2.2, 0.0) -- (2.2, 0.8);
\draw[->] (2.6, 0.0) -- (2.6, 0.8);

% Kurze Pfeile, die in das Innere des Ovals zeigen
\foreach \angle in {0, 45, 90, 135, 180, 225, 270, 315} {
    \draw[->, shorten >=1mm] (\angle:0.8cm and 0.6cm) ++(2.2,1.5) -- ++(\angle+180:0.3cm);
}

% ==============================
% 2. Kochtopf
% ==============================
\begin{scope}[xshift=6.5cm]
    \node at (-0.6, 3.425) [anchor=south west, inner sep=0] {Kochtopf};
    
    % Topfumriss
    \draw (-0.5, 3) -- (-0.5, 0.2) -- (1.2, 0.2) -- (1.2, 3);
    
    % Topfboden
    \draw (-0.7, 0) -- (1.4, 0);
    \draw (-0.6, 0) -- (-0.5, 0.2);
    \draw (1.3, 0) -- (1.2, 0.2);
    
    % Wasseroberfläche (wellige Linie)
    \draw[decorate, decoration={snake, amplitude=0.8mm, segment length=3mm}] (-0.5, 1.8) -- (1.2, 1.8);

    % Konvektionsströme (Zirkulationspfeile) im Wasser
    \draw[->] (-0.3, 0.5) to[out=0, in=180] (0.4, 0.8);  % unten links
    \draw[->] (0.9, 0.6) to[out=120, in=70] (0.5, 0.4); % unten rechts
    \draw[->] (-0.2, 1.3) to[out=90, in=120] (0.45, 1);  % oben links
    \draw[->] (0.9, 1.5) to[out=180, in=0] (0.4, 1.4);    % oben rechts
\end{scope}

% ==============================
% 3. Mikrowelle
% ==============================
\begin{scope}[xshift=9.25cm]
    \node at (1.15, 3.5) [anchor=south west, inner sep=0] {Mikrowelle};
    \draw (0,0) rectangle (4.5, 3);
    
    % Das Erhitzungsgut (Oval)
    \draw (2.25, 1.5) ellipse (1.2cm and 0.6cm);

    % Gekreuzte Pfeile (Dipolschwingungen) im Inneren
    \begin{scope}[shift={(1.5, 1.5)}]
        \draw[<->] (-0.3, 0) -- (0.3, 0);
        \draw[<->] (0, -0.3) -- (0, 0.3);
            \end{scope}
    \begin{scope}[shift={(3, 1.5)}]
        \draw[<->] (-0.3, 0) -- (0.3, 0);
        \draw[<->] (0, -0.3) -- (0, 0.3);
            \end{scope}
    \begin{scope}[shift={(2.25, 1.5)}]
        \draw[<->] (-0.3, 0) -- (0.3, 0);
        \draw[<->] (0, -0.3) -- (0, 0.3);
    \end{scope}
\end{scope}

\end{tikzpicture}
\caption{Ofen: Strahlung, Konvektion, Leitung ; Kochtopf: Konvektion, Leitung ; Mikrowelle: Erhitzung des Gutes von innen heraus}
\end{figure}




\end{document}
